\section{Folds: walking families}
\label{sec:walking}

The last rung of \S\ref{sec:taxonomy} is a count with no formula at all. This
section constructs the families that live there and the field they induce; the
next inverts them with a certificate. The tool that makes both tractable is
the same object the catalog's concentric rows were quietly built from.

\subsection{Radial metrics are support functions}

Fix a convex body $K\subset\R^2$ that is symmetric under rotation by $2\pi/N$ and
contains the origin. Write
\begin{equation}
  \rad(q) \;=\; \max_{u \in U}\, \langle q, u\rangle,
  \label{eq:support}
\end{equation}
where $U$ is the set of outward unit normals of $K$'s boundary. Then
$\rad(q)=R$ is precisely the boundary of the copy of $K$ with inradius $R$, so
$\rad$ is the natural ``radius'' for that shape. Equation~\eqref{eq:support} says
$\rad$ is the support function of $\mathrm{conv}(U)$~\cite{Rockafellar1970Convex}, and
the four cases our renderer exposes are all instances of it:
\begin{align}
  \text{circle:}\quad   &\rad(q) = \lVert q\rVert          &  U &= S^1, \notag\\
  \text{square:}\quad   &\rad(q) = \max(|q_x|, |q_y|)      &  U &= \{\pm e_1, \pm e_2\}, \notag\\
  \text{triangle:}\quad &\rad(q) = \max\Bigl(q_x,\; \tfrac{\sqrt3}{2}|q_y| - \tfrac{1}{2}q_x\Bigr), \notag\\
  \text{hexagon:}\quad  &\rad(q) = \max\Bigl(|q_x|,\; \tfrac{1}{2}|q_x| + \tfrac{\sqrt3}{2}|q_y|\Bigr).
  \label{eq:metrics}
\end{align}
The polygon forms in~\eqref{eq:metrics} are the max over $N$ unit normals written
out and folded by symmetry; for a free side count we evaluate
$\lVert q\rVert\cos(\mathrm{wrap}(\arg q, 2\pi/N))$ instead, which is the same
function. That the common cases avoid transcendentals entirely matters later:
$\rad$ is the innermost call of the solver.

We will use four properties, all immediate from~\eqref{eq:support}.

\begin{enumerate}
  \item[(P1)] \textbf{$1$-Lipschitz.} $|\rad(a)-\rad(b)| \le \lVert a-b\rVert$,
  since $U$ lies in the unit circle. This holds across corners, where $\nabla\rad$
  jumps.
  \item[(P2)] \textbf{Anisotropy floor.}
  $\kappa\lVert q\rVert \le \rad(q) \le \lVert q\rVert$ with
  \begin{equation}
    \kappa \;=\; \min_{\lVert v\rVert=1}\rad(v) \;=\; \cos(\pi/N),
    \label{eq:kappa}
  \end{equation}
  the inradius of $\mathrm{conv}(U)$: $1$ for the circle, $1/\sqrt2$ for the
  square, $1/2$ for the triangle. Both constants in~\eqref{eq:kappa} are attained,
  so neither can be improved.
  \item[(P3)] \textbf{Homogeneity and subadditivity.} $\rad(tq)=t\rad(q)$ for
  $t\ge0$ and $\rad(a+b)\le\rad(a)+\rad(b)$. Note that $\rad$ is \emph{not}
  symmetric: $\rad(-q)\neq\rad(q)$ for the triangle. This asymmetry is exactly
  what makes the drift constant of Section~\ref{sec:drift} directional.
  \item[(P4)] \textbf{Piecewise linearity.} For a polygon, $\rad$ is a max of $N$
  linear forms, hence convex and piecewise linear with $N$ pieces.
\end{enumerate}

\subsection{The family, and the field it induces}

\begin{figure}[t]
  \centering
  \begin{tikzpicture}[scale=0.92]
    % `\sp` is the inradius step, which is what the caption calls s. An equilateral
    % triangle's circumradius is twice its inradius; the factor here used to be
    % 1/cos(pi/6), which is the conversion for a *hexagon*, so the drawn triangles
    % were smaller than their stated inradius by 1/sqrt(3). The step is set so the
    % drawing is the size it always was, now under the right name.
    \def\sp{0.358}
    \def\dx{0.30}
    \def\dy{0.075}
    \def\th{13}
    % rings
    \foreach \n [evaluate=\n as \rr using 2*\n*\sp] [evaluate=\n as \cx using \n*\dx*cos(\n*\th) - \n*\dy*sin(\n*\th)] [evaluate=\n as \cy using \n*\dx*sin(\n*\th) + \n*\dy*cos(\n*\th)] in {1,2,3,4,5} {
      \draw[soft, thin, rotate around={\n*\th:(\cx,\cy)}]
        (\cx,\cy) +(90:\rr) -- +(210:\rr) -- +(330:\rr) -- cycle;
      \fill[soft] (\cx,\cy) circle (0.7pt);
    }
    % highlight ring 4
    \pgfmathsetmacro{\rr}{2*4*\sp}
    \pgfmathsetmacro{\cx}{4*\dx*cos(4*\th) - 4*\dy*sin(4*\th)}
    \pgfmathsetmacro{\cy}{4*\dx*sin(4*\th) + 4*\dy*cos(4*\th)}
    \draw[cool, thick, rotate around={4*\th:(\cx,\cy)}]
      (\cx,\cy) +(90:\rr) -- +(210:\rr) -- +(330:\rr) -- cycle;
    \fill[cool] (\cx,\cy) circle (1.4pt);
    \node[cool, font=\footnotesize, anchor=west] at (\cx+0.1,\cy-0.24) {$c_n$};
    % center walk
    \draw[accent, densely dotted, thick] (0,0)
      \foreach \n in {1,...,5} { -- ({\n*\dx*cos(\n*\th) - \n*\dy*sin(\n*\th)}, {\n*\dx*sin(\n*\th) + \n*\dy*cos(\n*\th)}) };
    \fill[ink] (0,0) circle (1.2pt);
    \node[ink, font=\footnotesize, anchor=north east] at (0,-0.05) {$O$};
    % query point. The arrow is drawn in world coordinates, so it is p - c_n; the
    % pulled-back point q_n is that same vector read in ring n's own frame, which
    % is a different vector, and rho is not rotation invariant -- labelling this
    % arrow q_n would undercut the very thing the figure is about.
    \coordinate (P) at (2.05,1.62);
    \fill[accent] (P) circle (1.6pt);
    \node[accent, font=\footnotesize, anchor=south west] at (2.1,1.66) {$p$};
    \draw[accent, thin, -{Stealth[length=4pt]}] (\cx,\cy) -- (P);
    \node[accent, font=\footnotesize] at (1.58,1.0) {$p-c_n$};
    \node[soft, font=\footnotesize, anchor=west] at (2.35,-1.6) {$n{=}1..5$};
  \end{tikzpicture}
  \Description{A diagram of five nested triangles of increasing size, each rotated a
  little further than the last and each centered a little further from the origin, so
  their centers trace a short dotted spiral outwards. The fourth triangle and its
  center are picked out in blue and labeled, and an arrow runs from that center to a
  query point marked inside the fourth triangle.}
  \caption{A walking family of triangles ($N=3$), drawn with $\varphi=0$. Ring $n$
  has inradius $ns+\varphi$, its center $c_n = \Rot{n\theta}(n\delta)$ walks along
  the dotted spiral, and its own frame is turned by $n\theta$. Rendering asks for
  the distance from $p$ to the nearest ring; the residual of ring $n$ is measured
  in that ring's frame, on the pulled-back point
  $q_n = \Rot{-n\theta}\,p - n\delta = \Rot{-n\theta}(p - c_n)$. The arrow is
  $p-c_n$, which is that same displacement in world coordinates: the two have equal
  length, but $\rad$ is not rotation invariant, so it is $q_n$ and not the arrow
  that the residual is taken of. That is the difficulty in one picture: because
  each ring has a different frame, no closed-form scalar field of $p$ has these
  rings as its level sets, and the index field survives only as the answer to a
  search (\S\ref{sec:inverse}).}
  \label{fig:family}
\end{figure}

A \emph{walking family} is the sequence of curves
\begin{equation}
  \Gamma_n \;=\; \bigl\{\, p : \rad\bigl(\Rot{-n\theta}(p - c_n)\bigr) = ns+\varphi \,\bigr\},
  \qquad c_n = \Rot{n\theta}(n\delta),
  \label{eq:family}
\end{equation}
for $n\in\N_0$, with spacing $s>0$, phase $\varphi\ge0$ (so that every radius
$ns+\varphi$ is nonnegative), per-index translation
$\delta\in\R^2$, and per-index rotation $\theta$. Figure~\ref{fig:family} shows
one. The three degenerate cases are familiar: $\delta=0,\theta=0$ gives
concentric shapes; $\theta=0$ alone gives a family whose centers walk a straight
line; and degenerating the shape to a half-plane, whose support set $U$ is a
single normal, recovers parallel lines at spacing $s$.

Expanding $\Rot{-n\theta}(p - \Rot{n\theta}(n\delta))$ gives the pulled-back
point
\begin{equation}
  q_n(p) \;=\; \Rot{-n\theta}\,p \;-\; n\delta,
  \label{eq:qn}
\end{equation}
which is the form we compute with: a rotation of $p$ that advances by $\theta$ per
index, minus a translation that advances by $\delta$ per index. Define the
\emph{index residual}
\begin{equation}
  h_p(n) \;=\; \rad\bigl(q_n(p)\bigr) - (ns + \varphi),
  \label{eq:residual}
\end{equation}
positive when $p$ falls outside ring $n$ and negative when inside. The distance
field the renderer wants is
\begin{equation}
  d(p) \;=\; \min_{n\in\N_0} \bigl|h_p(n)\bigr|.
  \label{eq:field}
\end{equation}
This is not the Euclidean distance to $\Gamma_{n^\star}$; it is the residual of
the inradius metric, which agrees with Euclidean distance for circles and
\emph{under}estimates it for polygons, by a factor of at most $1/\kappa$:
(P1) gives residual $\le$ distance always, and the gauge comparison gives
distance $\le$ residual$/\kappa$. Renderers of this kind consume $d$ through a
smoothstep of width one pixel, so the distinction thickens apparent stroke
weight near corners and changes nothing else. We keep the metric residual
because it is what admits the bounds below; the Limitations of
Section~\ref{sec:discussion} note the cost.

\subsection{The inversion problem}

Equation~\eqref{eq:field} is a minimization over an infinite discrete set, so
every practical solver is an argument about which finitely many $n$ can be
skipped. Three regimes are classical and one is not.

When $\delta=0$ and $\theta=0$, $h_p(n) = \rad(p)-ns-\varphi$ is affine in $n$,
so $d(p) = \mathrm{dist}\bigl((\rad(p)-\varphi) \bmod s\bigr)$ exactly: one
evaluation. When the shape is a circle and $\theta=0$, $\lVert p - n\delta\rVert =
ns+\varphi$ is a quadratic in $n$ whose roots bracket the answer; two integer
evaluations suffice. When $\theta=0$ and the shape is a polygon, $h_p$ is convex
and piecewise linear; we give the exact solution in
Section~\ref{sec:closed}. We have not found it stated elsewhere, and its
marginal case is the part that matters in practice.

Everything else (any $\theta\neq0$) has no closed form, because $q_n$ in
\eqref{eq:qn} contains $\Rot{-n\theta}p$, and asking for $h_p(n)=0$ asks where a
sinusoid in $n$ meets a line in $n$. That is transcendental, with a solution count
that depends on the parameters. The question is therefore not how to solve it, but
how few integers we can prove it suffices to test.

The naive answer is $n \approx \rad(p)/s$, ``the ring whose radius matches the
pixel's radius,'' plus a fixed number of neighbors. The next section shows how
far that can be from the truth, and replaces it with a provably sufficient
interval.
